% !TEX root = chapter-standalone.tex


\section{Querying}
\label{sec:dp-querying}

%\linkvideo{spring2021-functorial-comp-b:solving-queries} % Solving DP queries

A \SY{design problem} $\adp \colon \funspop \Ptimes \ressp \toinPos \Bool$ answers one question at a time: given a functionality $\fun$ and a requirement $\res$, is the pair $\tup{\fun, \res}$ feasible?
A designer, however, usually asks broader questions.
Typically, one of the two quantities is fixed and we ask about all possible values of the other:
\begin{itemize}
    \item Given a functionality $\fun$ that we must provide, which requirements $\res$ make $\tup{\fun, \res}$ feasible?
    \item Given a requirement $\res$, that is, required resources that are available to us, which functionalities $\fun$ can we provide with them?
\end{itemize}
For the engine of \cref{exa:engine-dp}, these questions are: ``how much fuel is enough to produce a given thrust?'' and ``what thrust can we produce with a given amount of fuel?''.
We call questions of this kind \emph{queries}.

In this section we show that the answers to queries are \SY{upper sets} and \SY{lower sets}, and that collecting all the answers to all queries gives an equivalent description of the design problem, obtained by \emph{currying} it.
At the end of the section we see how queries lead to optimization problems, which we study in \cref{sec:dp-optimality}.

\subsection{Queries and their answers}

\begin{definition}[Answers to queries]
    \label{def:dp-query-sets}
    Let $\adp \colon \funspop \Ptimes \ressp \toinPos \Bool$ be a \SY{design problem}.
    \begin{enumerate}
        \item For a functionality $\fun \setin \funsp$, the set of \emph{feasible requirements} for $\fun$ is
              \begin{equation}\label{eq:dp-feasible-requirements}
                  \widehat{\adp}(\fun) \definedas \makeset{ \res \setin \ressp \mid \adp(\fun\Fop, \res) = \true }.
              \end{equation}
        \item For a requirement $\res \setin \ressp$, the set of \emph{feasible functionalities} for $\res$ is
              \begin{equation}\label{eq:dp-feasible-functionalities}
                  \widecheck{\adp}(\res) \definedas \makeset{ \fun \setin \funsp \mid \adp(\fun\Fop, \res) = \true }.
              \end{equation}
    \end{enumerate}
\end{definition}

These sets are not arbitrary.
If the requirement $\res$ suffices to provide $\fun$, then so does any requirement $\res'$ with $\res \resleq \res'$, by the monotonicity of $\adp$ (this is condition \cref{eq:mon-rel-cond-1} for its feasible set).
And if $\fun$ can be provided with $\res$, then so can any functionality below $\fun$ (condition \cref{eq:mon-rel-cond-2}).

\begin{lemma}
    \label{lem:dp-query-sets}
    Let $\adp \colon \funspop \Ptimes \ressp \toinPos \Bool$ be a \SY{design problem}.
    For every $\fun \setin \funsp$, the set $\widehat{\adp}(\fun)$ is an \SY{upper set} of $\ressp$ (\cref{def:upperset}), and for every $\res \setin \ressp$, the set $\widecheck{\adp}(\res)$ is a \SY{lower set} of $\funsp$ (\cref{def:lowerset}).
\end{lemma}
% Proof commented out: it overlaps with the Fall 2026 homework (Sheet 2, CurryingDesignProblemsBasic).
% \begin{proof}
%     Let $\fun \setin \funsp$, and suppose $\res \setin \widehat{\adp}(\fun)$ and $\res \resleq \res'$.
%     Then $\tup{\fun\Fop, \res} \posleqof{\funspop \Ptimes \ressp} \tup{\fun\Fop, \res'}$, because $\fun\Fop \posleqof{\funspop} \fun\Fop$ by reflexivity.
%     Since $\adp$ is monotone, $\true = \adp(\fun\Fop, \res) \posleqof{\Bool} \adp(\fun\Fop, \res')$, and hence $\adp(\fun\Fop, \res') = \true$, that is, $\res' \setin \widehat{\adp}(\fun)$.
%
%     Now let $\res \setin \ressp$, and suppose $\fun \setin \widecheck{\adp}(\res)$ and $\fun' \funleq \fun$.
%     Then $\fun\Fop \posleqof{\funspop} {\fun'}\Fop$, so $\tup{\fun\Fop, \res} \posleqof{\funspop \Ptimes \ressp} \tup{{\fun'}\Fop, \res}$.
%     Since $\adp$ is monotone, $\adp({\fun'}\Fop, \res) = \true$, that is, $\fun' \setin \widecheck{\adp}(\res)$.
% \end{proof}

\begin{example}[The engine]
    \label{exa:engine-queries}
    For the engine of \cref{exa:engine-dp}, the set $\widehat{\text{engine}}(\fun)$ consists of all the amounts of fuel with which the thrust $\fun$ can be produced.
    Since $\Rtext{fuel}$ is totally ordered, an upper set of it is determined by a threshold: unless $\widehat{\text{engine}}(\fun)$ is empty (the thrust $\fun$ cannot be produced at all), it consists of all amounts of fuel from some minimal amount on.
    Similarly, $\widecheck{\text{engine}}(\res)$ consists of all the thrusts that can be produced with the amount of fuel $\res$, and it is a lower set of $\Ftext{thrust}$: all thrusts up to some maximal one (if the thrusts that can be produced with $\res$ have a maximum).
\end{example}

\subsection{Currying design problems}

How do the answers to queries change when we change the query?
Suppose that $\fun \funleq \fun'$, that is, $\fun'$ is a more demanding functionality than $\fun$.
If some requirement suffices to provide $\fun'$, then it also suffices to provide $\fun$.
So
\begin{equation}\label{eq:dp-queries-antitone}
    \prfperiod{
        \fun \funleq \fun'
    }{
        \widehat{\adp}(\fun') \setsubseteq \widehat{\adp}(\fun)
    }
\end{equation}
For a designer, this says that a more demanding functionality leaves fewer feasible requirements.
Dually, if $\res \resleq \res'$, then $\widecheck{\adp}(\res) \setsubseteq \widecheck{\adp}(\res')$: a larger requirement, that is, more of the required resources, gives more choices of functionality.

The reversal in \cref{eq:dp-queries-antitone} is exactly what the opposite order on the functionalities in the definition of a design problem accounts for.
Write $\tup{\uppersets \ressp, \setsubseteq}$ for the set of \SY{upper sets} of $\ressp$, ordered by inclusion, and $\tup{\lowersets \funsp, \setsubseteq}$ for the set of \SY{lower sets} of $\funsp$, ordered by inclusion.

\begin{lemma}
    \label{lem:dp-currying-monotone}
    Let $\adp \colon \funspop \Ptimes \ressp \toinPos \Bool$ be a \SY{design problem}.
    Then the functions
    \begin{equation}\label{eq:dp-left-curry}
        \defmapcomma{
            \widehat{\adp}
        }{
            \funspop
        }{
            \mto
        }{
            \tup{\uppersets \ressp, \setsubseteq}
        }{
            \fun\Fop
        }{
            \widehat{\adp}(\fun)
        }
    \end{equation}
    and
    \begin{equation}\label{eq:dp-right-curry}
        \defmapperiod{
            \widecheck{\adp}
        }{
            \ressp
        }{
            \mto
        }{
            \tup{\lowersets \funsp, \setsubseteq}
        }{
            \res
        }{
            \widecheck{\adp}(\res)
        }
    \end{equation}
    are \SY{monotone functions}.
\end{lemma}
% Proof commented out: it overlaps with the Fall 2026 homework (Sheet 2, CurryingDesignProblemsBasic).
% \begin{proof}
%     By \cref{lem:dp-query-sets}, both functions take values in the stated sets.
%     For $\widehat{\adp}$: suppose $\fun\Fop \posleqof{\funspop} {\fun'}\Fop$, that is, $\fun' \funleq \fun$, and let $\res \setin \widehat{\adp}(\fun)$.
%     Then $\tup{\fun\Fop, \res} \posleqof{\funspop \Ptimes \ressp} \tup{{\fun'}\Fop, \res}$, so, since $\adp$ is monotone, $\adp({\fun'}\Fop, \res) = \true$, that is, $\res \setin \widehat{\adp}(\fun')$.
%     This shows $\widehat{\adp}(\fun) \setsubseteq \widehat{\adp}(\fun')$.
%     For $\widecheck{\adp}$: suppose $\res \resleq \res'$, and let $\fun \setin \widecheck{\adp}(\res)$.
%     Then $\tup{\fun\Fop, \res} \posleqof{\funspop \Ptimes \ressp} \tup{\fun\Fop, \res'}$, so $\adp(\fun\Fop, \res') = \true$, that is, $\fun \setin \widecheck{\adp}(\res')$.
%     This shows $\widecheck{\adp}(\res) \setsubseteq \widecheck{\adp}(\res')$.
% \end{proof}

\begin{remark}
    \label{rem:dp-curry-other-form}
    A monotone function $\posA\posop \mto \posB$ is the same thing as a monotone function $\posA \mto \posB\posop$, with the same underlying function.
    Since the opposite of $\tup{\uppersets \ressp, \setsubseteq}$ is $\tup{\uppersets \ressp, \setsupseteq}$, the function $\widehat{\adp}$ of \cref{eq:dp-left-curry} can equivalently be seen as a monotone function
    \begin{equation}
        \widehat{\adp} \colon \funsp \mto \tup{\uppersets \ressp, \setsupseteq},
    \end{equation}
    from the functionalities, with their original order, to the upper sets ordered by \emph{reverse} inclusion.
    This is just \cref{eq:dp-queries-antitone} again.
\end{remark}

\subsubsection{Why this is currying}

The name ``currying'' comes from the construction of \cref{sec:arithmetic-exponents-products}.
There, a function of two variables $\mora \colon \setA \cartprod \setB \sto \setC$ was turned into the function $\curry(\mora) \colon \setA \sto \setC^\setB$, which sends $\ela$ to the function $\mora(\ela, -)$ of the remaining variable (\cref{eq:left-curry-function}).
We now explain why $\widehat{\adp}$ is the analogue of this construction for design problems, and why this makes the order ``$\setsubseteq$'' on upper sets the natural choice.

Currying a design problem $\adp \colon \funspop \Ptimes \ressp \toinPos \Bool$ in the same way, we obtain, for each $\fun\Fop \setin \funspop$, the function
\begin{equation}
    \defmapperiod{
        \adp(\fun\Fop, -)
    }{
        \ressp
    }{
        \mto
    }{
        \Bool
    }{
        \res
    }{
        \adp(\fun\Fop, \res)
    }
\end{equation}
It is monotone, since $\adp$ is monotone and $\fun\Fop$ is kept fixed.
So currying lands in the set of monotone functions $\ressp \mto \Bool$.
The following lemma identifies this set, together with its natural order, with the upper sets of $\ressp$.
Here we order monotone functions $\phi, \psi \colon \ressp \mto \Bool$ point-wise: $\phi \posleq \psi$ if and only if $\phi(\res) \posleqof{\Bool} \psi(\res)$ for all $\res \setin \ressp$.

\begin{lemma}
    \label{lem:upper-sets-as-maps-to-bool}
    Let $\posA$ be a \SY{poset}.
    \begin{enumerate}
        \item For every monotone function $\phi \colon \posA \mto \Bool$, the set $\stylesets{U}_\phi = \makeset{\posAel \setin \posAset \mid \phi(\posAel) = \true}$ is an \SY{upper set} of $\posA$, and every upper set of $\posA$ is of this form for exactly one monotone function $\phi$.
        \item For monotone functions $\phi, \psi \colon \posA \mto \Bool$, we have $\phi \posleq \psi$ point-wise if and only if $\stylesets{U}_\phi \setsubseteq \stylesets{U}_\psi$.
    \end{enumerate}
\end{lemma}
\begin{proof}
    % Proof of part 1 commented out: it overlaps with the Fall 2026 homework (Sheet 2, UpperSetsAsMonotoneFunctions).
    % For part 1: if $\phi$ is monotone, $\phi(\posAel) = \true$, and $\posAel \posAleq \posAel'$, then $\true = \phi(\posAel) \posleqof{\Bool} \phi(\posAel')$, so $\phi(\posAel') = \true$; hence $\stylesets{U}_\phi$ is an upper set.
    % Conversely, given an upper set $\stylesets{U}$, its indicator function $\phi$ (with $\phi(\posAel) = \true$ if and only if $\posAel \setin \stylesets{U}$) is monotone: if $\posAel \posAleq \posAel'$, then either $\phi(\posAel) = \false$, and $\false \posleqof{\Bool} \phi(\posAel')$, or $\posAel \setin \stylesets{U}$, and then $\posAel' \setin \stylesets{U}$, so $\phi(\posAel') = \true$.
    % A function to $\boolset$ is determined by the set of elements it sends to $\true$, so $\phi$ is the only function with $\stylesets{U}_\phi = \stylesets{U}$.
    We prove part 2.
    Suppose first that $\phi \posleq \psi$ point-wise.
    If $\posAel \setin \stylesets{U}_\phi$, then $\true = \phi(\posAel) \posleqof{\Bool} \psi(\posAel)$, so $\psi(\posAel) = \true$, that is, $\posAel \setin \stylesets{U}_\psi$.
    Conversely, suppose $\stylesets{U}_\phi \setsubseteq \stylesets{U}_\psi$, and let $\posAel \setin \posAset$.
    If $\phi(\posAel) = \false$, then $\phi(\posAel) \posleqof{\Bool} \psi(\posAel)$ because $\false$ is the least element of $\Bool$.
    If $\phi(\posAel) = \true$, then $\posAel \setin \stylesets{U}_\phi \setsubseteq \stylesets{U}_\psi$, so $\psi(\posAel) = \true$, and again $\phi(\posAel) \posleqof{\Bool} \psi(\posAel)$.
\end{proof}

Under the correspondence of \cref{lem:upper-sets-as-maps-to-bool}, the function $\adp(\fun\Fop, -)$ corresponds to the upper set $\makeset{\res \mid \adp(\fun\Fop, \res) = \true}$, which is $\widehat{\adp}(\fun)$.
So in this sense $\widehat{\adp}$ is the curried form of $\adp$, with monotone functions $\ressp \mto \Bool$ replaced by the corresponding upper sets.
And the point-wise order on monotone functions becomes inclusion of upper sets.
In this sense the order ``$\setsubseteq$'' in \cref{eq:dp-left-curry} is not an arbitrary choice: it is the order that currying produces.

This is also the analogue, for \SY{posets}, of the description of relations in \cref{rem:rel-three-fun-descriptions}.
A relation $\relA \setsubseteq \setA \cartprod \setB$ can be described by the function $\hat\phi_\relA \colon \setA \sto \powerset(\setB)$, and subsets of $\setB$ correspond to functions $\setB \sto \boolset$.
For a design problem, subsets are replaced by upper sets, which correspond to \emph{monotone} functions $\ressp \mto \Bool$.

As for relations, no information is lost in passing to the curried form.

\begin{proposition}
    \label{prop:dp-currying}
    Let $\funsp$ and $\ressp$ be \SY{posets}.
    Sending a design problem $\adp$ to $\widehat{\adp}$ gives a one-to-one correspondence between \SY{design problems} $\adp \colon \funspop \Ptimes \ressp \toinPos \Bool$ and monotone functions $\funspop \mto \tup{\uppersets \ressp, \setsubseteq}$.
    Its inverse sends a monotone function $\mapc \colon \funspop \mto \tup{\uppersets \ressp, \setsubseteq}$ to the design problem $\adp_\mapc$ given by
    \begin{equation}
        \adp_\mapc(\fun\Fop, \res) = \true
        \quad \text{if and only if} \quad
        \res \setin \mapc(\fun\Fop).
    \end{equation}
    Similarly, sending $\adp$ to $\widecheck{\adp}$ gives a one-to-one correspondence between design problems and monotone functions $\ressp \mto \tup{\lowersets \funsp, \setsubseteq}$.
\end{proposition}
% Proof commented out: the monotonicity checks overlap with the Fall 2026 homework (Sheet 2, CurryingDesignProblemsBasic).
% \begin{proof}
%     We treat $\widehat{\adp}$; the case of $\widecheck{\adp}$ is analogous.
%     By \cref{lem:dp-currying-monotone}, $\widehat{\adp}$ is a monotone function of the stated type.
%     Conversely, let $\mapc$ be such a monotone function; we check that $\adp_\mapc$ is monotone.
%     Suppose $\tup{\fun\Fop, \res} \posleqof{\funspop \Ptimes \ressp} \tup{{\fun'}\Fop, \res'}$, that is, $\fun\Fop \posleqof{\funspop} {\fun'}\Fop$ and $\res \resleq \res'$, and suppose $\adp_\mapc(\fun\Fop, \res) = \true$, that is, $\res \setin \mapc(\fun\Fop)$.
%     Since $\mapc(\fun\Fop)$ is an upper set, $\res' \setin \mapc(\fun\Fop)$; since $\mapc$ is monotone, $\mapc(\fun\Fop) \setsubseteq \mapc({\fun'}\Fop)$, so $\res' \setin \mapc({\fun'}\Fop)$, that is, $\adp_\mapc({\fun'}\Fop, \res') = \true$.
%     As in the proof of \cref{lem:upper-sets-as-maps-to-bool}, a function to $\Bool$ is monotone as soon as $\true$ values propagate upward, so $\adp_\mapc$ is monotone.
%     Finally, the two constructions are inverse to one another: $\res \setin \widehat{\adp_\mapc}(\fun)$ if and only if $\res \setin \mapc(\fun\Fop)$, and $\adp_{\widehat{\adp}}(\fun\Fop, \res) = \true$ if and only if $\res \setin \widehat{\adp}(\fun)$, that is, if and only if $\adp(\fun\Fop, \res) = \true$.
% \end{proof}

So we have three equivalent descriptions of a design problem: the monotone function $\adp$ itself, its left curried form $\widehat{\adp}$, which answers all queries of the form ``given $\fun$, which requirements?'', and its right curried form $\widecheck{\adp}$, which answers all queries of the form ``given $\res$, which functionalities?''.

\begin{example}[A finite design problem]
    \label{exa:finite-dp-queries}
    Consider the design problem $\adp$ of \cref{exa:visualize-dp} (\cref{fig:example_dp_graph_xy}), in which $\setAeln{1} \posleq \setAeln{2} \posleq \setAeln{3}$ and $\setBeln{1} \posleq \setBeln{2}$.
    Reading off the feasible pairs, we find
    \begin{equation}
        \widehat{\adp}(\setAeln{1}) = \makeset{\setBeln{1}, \setBeln{2}},
        \qquad
        \widehat{\adp}(\setAeln{2}) = \makeset{\setBeln{2}},
        \qquad
        \widehat{\adp}(\setAeln{3}) = \makeset{\setBeln{2}},
    \end{equation}
    and
    \begin{equation}
        \widecheck{\adp}(\setBeln{1}) = \makeset{\setAeln{1}},
        \qquad
        \widecheck{\adp}(\setBeln{2}) = \makeset{\setAeln{1}, \setAeln{2}, \setAeln{3}}.
    \end{equation}
    Each $\widehat{\adp}(\setAeln{i})$ is an upper set of $\setBeln{1} \posleq \setBeln{2}$, and each $\widecheck{\adp}(\setBeln{j})$ is a lower set of $\setAeln{1} \posleq \setAeln{2} \posleq \setAeln{3}$.
    Moreover, going up in the functionalities makes the sets of feasible requirements shrink, $\widehat{\adp}(\setAeln{1}) \setsupseteq \widehat{\adp}(\setAeln{2}) \setsupseteq \widehat{\adp}(\setAeln{3})$, while going up in the requirements makes the sets of feasible functionalities grow, $\widecheck{\adp}(\setBeln{1}) \setsubseteq \widecheck{\adp}(\setBeln{2})$.
\end{example}

\begin{example}[Design problems from monotone functions]
    \label{exa:companion-dp-queries}
    Let $\mapb \colon \funsp \toinPos \ressp$ be a monotone function, and let $\adp_\mapb$ be the design problem of \cref{ex:DPsFromMonotoneMaps}, with $\adp_\mapb(\fun\Fop, \res) = \true$ if and only if $\mapb(\fun) \resleq \res$.
    We can think of $\mapb(\fun)$ as ``the requirement of $\fun$''.
    Then
    \begin{equation}
        \widehat{\adp_\mapb}(\fun) = \makeset{\res \setin \ressp \mid \mapb(\fun) \resleq \res},
    \end{equation}
    which is the upper set of all requirements above the single element $\mapb(\fun)$.
    So every query ``given $\fun$, which requirements?'' has an answer with a least element, namely $\mapb(\fun)$.
    On the other hand,
    \begin{equation}
        \widecheck{\adp_\mapb}(\res) = \makeset{\fun \setin \funsp \mid \mapb(\fun) \resleq \res}
    \end{equation}
    is the lower set of all functionalities whose requirement $\mapb(\fun)$ is at most $\res$; in general it has no greatest element.
\end{example}

\begin{example}[Paying with coins]
    \label{exa:coins-dp}
    Let $\funsp = \natswithleq$, whose elements are amounts of money (in francs) that we need to pay, and let $\ressp = \natswithleq \Ptimes \natswithleq$, whose elements $\tup{\elnc{1}, \elnc{2}}$ are numbers of one-franc and two-franc coins, ordered component-wise.
    The design problem
    \begin{equation}
        \adp(\ela\Fop, \tup{\elnc{1}, \elnc{2}}) = \true
        \quad \text{if and only if} \quad
        \ela \leq \elnc{1} + 2\elnc{2}
    \end{equation}
    says that the coins suffice to pay the amount (possibly overpaying).
    For the amount $3$, the set of feasible requirements is
    \begin{equation}
        \widehat{\adp}(3) = \makeset{\tup{\elnc{1}, \elnc{2}} \mid 3 \leq \elnc{1} + 2\elnc{2}}.
    \end{equation}
    It contains $\tup{3, 0}$, $\tup{1, 1}$, and $\tup{0, 2}$, and each of these three combinations is \emph{minimal}: removing any coin from it leaves less than three francs.
    They are pairwise incomparable, so, unlike in \cref{exa:companion-dp-queries}, $\widehat{\adp}(3)$ has no least element: there is no single ``cheapest'' way to pay three francs.
    In the other direction, $\widecheck{\adp}(\tup{\elnc{1}, \elnc{2}}) = \makeset{0, 1, \dots, \elnc{1} + 2\elnc{2}}$: with given coins we can pay any amount up to their total value.
\end{example}

\subsection{From queries to optimization}

In practice, a designer rarely wants the whole set $\widehat{\adp}(\fun)$ of feasible requirements.
Since it is an upper set, it contains, together with any feasible requirement, all the more wasteful requirements above it.
What we really want to know are the \emph{minimal} feasible requirements, that is, the set $\Min \widehat{\adp}(\fun)$ of minimal elements (\cref{def:Min}): the requirements that cannot be lowered without becoming infeasible.
Dually, given a requirement $\res$, we want to know the \emph{maximal} feasible functionalities $\Max \widecheck{\adp}(\res)$ (\cref{def:Max}).
Later in the book, we will call these two questions $\FixFunMinRes$ (``fix functionality, minimize the required resources'') and $\FixResMaxFun$ (``fix the required resources, maximize functionality'').

The examples above show what kind of answers to expect.
For the engine (\cref{exa:engine-queries}) and for design problems coming from monotone functions (\cref{exa:companion-dp-queries}), the minimal feasible requirements consist of at most one element.
For the coins (\cref{exa:coins-dp}), we have
\begin{equation}
    \Min \widehat{\adp}(3) = \makeset{\tup{3, 0}, \tup{1, 1}, \tup{0, 2}},
\end{equation}
a set of pairwise incomparable elements, that is, an \SY{antichain}.
This is typical: when requirements are only partially ordered, there is in general no single best choice, but a set of incomparable optimal trade-offs.
In \cref{sec:dp-optimality} we study \SY{antichains} and their relation to \SY{upper sets} in detail.


% The former graded exercise, kept commented out:
%\begin{gradedexercise}[\exname{CurryingDesignProblems}]
%    \label{ex:CurryingDesignProblems}
%    Let $\adp\colon \funspop \Ptimes \ressp\toinPos \Bool$ be a design problem.
%    In this exercise we will show that \cref{eq:dp-left-curry} corresponds to a monotone function
%    \begin{equation}
%        \funsp \mto \tupp{\uppersets (\ressp), {{\setsupseteq}}},
%    \end{equation}
%    and that \cref{eq:dp-right-curry} corresponds to a monotone function
%    \begin{equation}
%        \ressp \mto \tupp{\setOfLowersets(\funsp), \setsubseteq}.
%    \end{equation}
%    Here $\uppersets (\ressp)$ denotes the set of \SY{upper sets} of $\ressp$ and $\setOfLowersets(\funsp)$ denotes the set of \SY{lower sets} of $\funsp$.
%
%    \begin{enumerate}
%        \item Show that $\widehat{\adp} \colon \funspop \mto \uppersets (\ressp)$ and $\widecheck{\adp} \colon \ressp \mto \uppersets (\funspop)$ are \SY{monotone maps} when we consider $\uppersets (\ressp)$ and $\uppersets (\funspop)$ to have the partial order corresponding to the inclusion of subsets.
%        \item Show that the \SY{poset} $\tupp{\uppersets\funspop,\setsubseteq}$ and the \SY{poset} $\tupp{\setOfLowersets\funsp, \setsubseteq}$ are isomorphic.
%        \item Show that there is a 1-to-1 correspondence between  \SY{monotone functions} \begin{equation}
%                  \funspop \mto \tupp{\uppersets \ressp, \setsubseteq}
%              \end{equation}
%              and  \SY{monotone functions} \begin{equation}
%                  \funsp \mto \tupp{\uppersets \ressp, {{\setsupseteq}}},
%              \end{equation}
%              where in the latter poset, the order is given by the relation of ``containment'' (as opposed to ``inclusion'').
%        \item Explain, in a few words, why the above steps prove the stated goal of this exercise.
%    \end{enumerate}
%\end{gradedexercise}
%\solutionof{CurryingDesignProblems}
